\documentclass[12pt,letterpaper]{article}

% ------------------------------------------------
% Paquetes
% ------------------------------------------------
\usepackage[utf8]{inputenc}
\usepackage[T1]{fontenc}
\usepackage[spanish]{babel}

\usepackage{amsmath,amssymb,amsthm,mathtools,bm}
\usepackage{booktabs}
\usepackage{array}
\usepackage{geometry}
\usepackage{setspace}
\usepackage{microtype}
\usepackage{parskip}
\usepackage{titlesec}
\usepackage{fancyhdr}
\usepackage{enumitem}
\usepackage{xcolor}
\usepackage{tcolorbox}
\usepackage{hyperref}

\tcbuselibrary{skins,breakable}

% ------------------------------------------------
% Geometría
% ------------------------------------------------
\geometry{
  top=2.8cm,
  bottom=3cm,
  left=3.2cm,
  right=3.2cm,
  headheight=14pt
}

\setstretch{1.20}

% ------------------------------------------------
% Encabezados
% ------------------------------------------------
\pagestyle{fancy}
\fancyhf{}
\fancyhead[C]{\small\textsc{Econometría --- Gujarati --- Ejercicio 13.10}}
\fancyfoot[C]{\small\thepage}

\renewcommand{\headrulewidth}{0.4pt}

% ------------------------------------------------
% Títulos
% ------------------------------------------------
\titleformat{\section}
{\large\bfseries\scshape}
{\thesection.}{0.6em}{}
[\vspace{-4pt}\rule{\linewidth}{0.4pt}]

\titleformat{\subsection}
{\normalsize\bfseries}
{\thesubsection.}{0.5em}{}

\titleformat{\subsubsection}[runin]
{\normalsize\bfseries}
{}{0em}{}[.\quad]

% ------------------------------------------------
% Caja destacada
% ------------------------------------------------
\newtcolorbox{clave}{
  colback=white,
  colframe=black,
  boxrule=0.5pt,
  arc=3pt,
  left=8pt,
  right=8pt,
  top=6pt,
  bottom=6pt,
  breakable
}

% ------------------------------------------------
% Comandos útiles
% ------------------------------------------------
\newcommand{\Var}{\operatorname{Var}}
\newcommand{\Cov}{\operatorname{Cov}}
\newcommand{\E}{\operatorname{E}}
\newcommand{\plim}{\operatorname{plim}}
\newcommand{\RSS}{\operatorname{RSS}}

% ==========================================================
\begin{document}
% ==========================================================

\begin{center}

{\LARGE\bfseries Ejercicio 13.10 --- Prueba del Multiplicador}\\[4pt]

{\LARGE\bfseries de Lagrange y Variable Omitida}\\[12pt]

{\large\itshape Econometría --- Gujarati \& Porter}\\[10pt]

\rule{0.55\linewidth}{0.6pt}\\[10pt]

{\normalsize
Ejercicio 13.10 resuelto por Emanuel Quintana Silva}\\[4pt]

{\small
Regresión auxiliar, prueba LM y por qué incluir todos los regresores}

\end{center}

\vspace{12pt}

% ==========================================================
\section{Enunciado}
% ==========================================================

Considere el modelo:

\begin{equation}
Y_{i} = \beta_{1} + \beta_{2} X_{2i} + u_{i}
\end{equation}

Para averiguar si este modelo está mal especificado por omitir la
variable $X_{3}$, se propone regresar los residuos $\hat{u}_{i}$
obtenidos de (1) únicamente sobre $X_{3}$ (incluyendo un intercepto).

La prueba del \textbf{Multiplicador de Lagrange (ML)}, sin embargo,
requiere regresar los residuos de (1) sobre $X_{2}$, $X_{3}$ y una
constante.

\bigskip

\noindent\textit{¿Por qué es probable que el procedimiento propuesto
sea inapropiado?}

% ==========================================================
\section{Contexto: la prueba del Multiplicador de Lagrange}
% ==========================================================

La prueba LM para detectar variables omitidas sigue el siguiente
procedimiento estándar:

\begin{enumerate}[leftmargin=1.8em]
\item Estimar el \textbf{modelo restringido} (1) por MCO y obtener los
      residuos $\hat{u}_{i}$.
\item Estimar la \textbf{regresión auxiliar}:
      \begin{equation}
        \hat{u}_{i} = \gamma_{0} + \gamma_{1} X_{2i} + \gamma_{2} X_{3i}
                      + \varepsilon_{i}
      \end{equation}
      incluyendo \emph{todos} los regresores: tanto los del modelo
      original como la variable sospechosa de haber sido omitida.
\item Calcular el estadístico:
      \begin{equation}
        LM = n\, R^{2}_{\hat{u}}
      \end{equation}
      donde $R^{2}_{\hat{u}}$ es el coeficiente de determinación de la
      regresión (2) y $n$ es el tamaño de la muestra.
\item Bajo $H_{0}$: $\beta_{3} = 0$, el estadístico $LM$ sigue
      asintóticamente una distribución:
      \begin{equation}
        LM \xrightarrow{d} \chi^{2}(q)
      \end{equation}
      donde $q$ es el número de restricciones impuestas (en este caso,
      $q = 1$).
\end{enumerate}

El procedimiento \emph{propuesto} omite $X_{2}$ de la regresión
auxiliar, estimando en su lugar:

\begin{equation}
\hat{u}_{i} = \delta_{0} + \delta_{1} X_{3i} + \xi_{i}
\end{equation}

A continuación se demuestra por qué esta simplificación es inválida.

% ==========================================================
\section{Razón 1: Correlación entre $X_{2}$ y $X_{3}$}
% ==========================================================

\subsection{Sesgo por variable omitida en la primera etapa}

Si $X_{3}$ es una variable relevante \emph{y} está correlacionada
con $X_{2}$, es decir, si:

\begin{equation}
\Cov(X_{2i},\, X_{3i}) \neq 0
\end{equation}

entonces el estimador MCO del modelo restringido (1) satisface:

\begin{equation}
\plim(\hat{\beta}_{2})
= \beta_{2} + \beta_{3}\,\delta_{32}
\end{equation}

donde $\delta_{32}$ es el coeficiente de la regresión auxiliar de
$X_{3}$ sobre $X_{2}$:

\begin{equation}
\delta_{32}
= \frac{\Cov(X_{2i},\, X_{3i})}{\Var(X_{2i})}
\end{equation}

Por lo tanto, $\hat{\beta}_{2}$ es \textbf{sesgado e inconsistente}:
absorbe parte de la influencia de $X_{3}$ sobre $Y$.

\subsection{Efecto sobre los residuos $\hat{u}_{i}$}

Los residuos del modelo restringido se escriben como:

\begin{equation}
\hat{u}_{i}
= Y_{i} - \hat{\beta}_{1} - \hat{\beta}_{2} X_{2i}
\end{equation}

Por construcción de MCO, estos residuos son \textbf{ortogonales} a
$X_{2}$:

\begin{equation}
\sum_{i=1}^{n} X_{2i}\,\hat{u}_{i} = 0
\end{equation}

Sin embargo, todavía contienen la influencia de $X_{3}$ que no fue
capturada correctamente, ya que $\hat{\beta}_{2}$ absorbió parte de
ella. El efecto de $X_{3}$ sobre $\hat{u}_{i}$ no puede separarse
del efecto de $X_{2}$ sin incluir a $X_{2}$ en la regresión auxiliar.

% ==========================================================
\section{Razón 2: Requisito de la regresión auxiliar completa}
% ==========================================================

Al incluir únicamente $X_{3}$ en la regresión auxiliar (ecuación 4),
el estimador $\hat{\delta}_{1}$ satisface:

\begin{equation}
\plim(\hat{\delta}_{1})
= \gamma_{2} + \gamma_{1}\,\frac{\Cov(X_{2i},\, X_{3i})}{\Var(X_{3i})}
\end{equation}

Es decir, $\hat{\delta}_{1}$ confunde la contribución marginal de
$X_{3}$ con la influencia que $X_{2}$ tiene sobre los residuos a
través de su correlación con $X_{3}$.

Solo cuando se incluyen \emph{todos} los regresores en la ecuación
auxiliar (ecuación 2), los coeficientes $\gamma_{1}$ y $\gamma_{2}$
estiman correctamente los efectos parciales, y el estadístico $LM$
queda bien definido.

\begin{clave}
Wooldridge señala explícitamente que en la regresión auxiliar de la
prueba LM \textbf{se deben incluir todos los regresores} porque, en
general, los regresores omitidos en el modelo restringido están
correlacionados con los que sí aparecen en él. Omitir $X_{2}$ impide
aislar correctamente la contribución marginal de $X_{3}$.
\end{clave}

% ==========================================================
\section{Razón 3: Invalidez del estadístico de prueba}
% ==========================================================

\subsection{Distribución del estadístico bajo $H_{0}$}

La validez asintótica del estadístico $LM = n R^{2}_{\hat{u}}$ depende
de que la regresión auxiliar incluya el conjunto completo de regresores.
Solo en ese caso se puede demostrar que:

\begin{equation}
LM = n\,R^{2}_{\hat{u}} \xrightarrow{d} \chi^{2}(q)
\quad \text{bajo } H_{0}
\end{equation}

\subsection{Qué ocurre con la regresión ``corta''}

Sea $\tilde{R}^{2}$ el coeficiente de determinación de la regresión
auxiliar incompleta (ecuación 4). En general:

\begin{equation}
\tilde{R}^{2} \neq R^{2}_{\hat{u}}
\end{equation}

El estadístico $n\tilde{R}^{2}$ \textbf{no sigue} la distribución
$\chi^{2}(1)$ bajo $H_{0}$. Por lo tanto, cualquier conclusión
inferencial basada en él es inválida: podría llevar tanto a
\emph{rechazar} erróneamente un modelo correctamente especificado
como a \emph{no rechazar} uno mal especificado.

% ==========================================================
\section{Demostración algebraica de la invalidez}
% ==========================================================

Partimos del modelo verdadero:

\begin{equation}
Y_{i} = \beta_{1} + \beta_{2} X_{2i} + \beta_{3} X_{3i} + \varepsilon_{i}
\end{equation}

Los residuos del modelo restringido (1) son:

\begin{align}
\hat{u}_{i}
&= Y_{i} - \hat{\beta}_{1} - \hat{\beta}_{2} X_{2i} \\[4pt]
&= \beta_{3} X_{3i} + \varepsilon_{i}
   + (\beta_{2} - \hat{\beta}_{2}) X_{2i}
   + (\beta_{1} - \hat{\beta}_{1})
\end{align}

Si se regresa $\hat{u}_{i}$ únicamente sobre $X_{3}$, el estimador
$\hat{\delta}_{1}$ recoge tanto el efecto de $\beta_{3}$ como el
sesgo $(\beta_{2} - \hat{\beta}_{2})$ transmitido a través de
$\Cov(X_{2i}, X_{3i})$. En cambio, si se incluye $X_{2}$ en la
regresión auxiliar, la ortogonalidad de los residuos MCO respecto
a $X_{2}$ garantiza que:

\begin{equation}
\sum_{i=1}^{n} X_{2i}\,\hat{u}_{i} = 0
\end{equation}

y el coeficiente de $X_{3}$ en la regresión auxiliar completa
identifica \emph{exclusivamente} el efecto marginal de $X_{3}$ no
capturado por el modelo restringido.

% ==========================================================
\section{Resumen del argumento}
% ==========================================================

\begin{center}
\renewcommand{\arraystretch}{1.5}
\begin{tabular}{p{4.5cm}p{4.8cm}p{4.8cm}}
\toprule
\textbf{Aspecto}
  & \textbf{Procedimiento propuesto}
  & \textbf{Prueba LM correcta} \\
\midrule
Regresión auxiliar
  & $\hat{u}_{i}$ sobre $X_{3}$ solo
  & $\hat{u}_{i}$ sobre $X_{2}$, $X_{3}$ y constante \\
Efecto de $\Cov(X_{2},X_{3})$
  & Confunde efectos de $X_{2}$ y $X_{3}$
  & Aísla el efecto marginal de $X_{3}$ \\
$R^{2}$ de la auxiliar
  & $\tilde{R}^{2}$ no tiene distribución conocida
  & $R^{2}_{\hat{u}}$ sigue $\chi^{2}(q)/n$ bajo $H_{0}$ \\
Estadístico $n R^{2}$
  & \textbf{Inválido}: distribución desconocida
  & \textbf{Válido}: $\chi^{2}(1)$ bajo $H_{0}$ \\
Conclusión inferencial
  & Potencialmente errónea
  & Asintóticamente correcta \\
\bottomrule
\end{tabular}
\end{center}

% ==========================================================
\section{Conclusión}
% ==========================================================

\begin{clave}

El procedimiento propuesto es inapropiado por tres razones
interdependientes:

\begin{enumerate}[leftmargin=1.8em]

\item \textbf{Sesgo en la primera etapa.} Si $\Cov(X_{2i}, X_{3i})
      \neq 0$, los residuos $\hat{u}_{i}$ del modelo restringido
      contienen el efecto de $X_{3}$ mezclado con el sesgo de
      $\hat{\beta}_{2}$. Solo incluyendo $X_{2}$ en la regresión
      auxiliar se purga este efecto cruzado.

\item \textbf{Estimación parcial inválida.} Regresar $\hat{u}_{i}$
      solo sobre $X_{3}$ no aísla la contribución marginal de
      $X_{3}$: el coeficiente estimado confunde el efecto directo
      de $X_{3}$ con su correlación con $X_{2}$.

\item \textbf{Estadístico sin distribución conocida.} El estadístico
      $n\tilde{R}^{2}$ de la regresión ``corta'' no sigue la
      distribución $\chi^{2}(1)$ bajo $H_{0}$, por lo que cualquier
      inferencia basada en él carece de sustento asintótico.

\end{enumerate}

La prueba LM es válida únicamente cuando la regresión auxiliar incluye
el \textbf{conjunto completo de regresores}: los del modelo original
más la variable cuya omisión se investiga.

\end{clave}

\vspace{12pt}

\noindent\rule{\linewidth}{0.4pt}

\vspace{10pt}

\begin{center}
\small\textsc{Referencias}
\end{center}

\vspace{6pt}

\noindent
Gujarati, D.~N., \& Porter, D.~C. (2010).
\textit{Econometría} (5.a ed.).
McGraw-Hill.

\vspace{6pt}

\noindent
Wooldridge, J.~M. (2016).
\textit{Introductory Econometrics} (6.a ed.).
Cengage.

\vspace{6pt}

\noindent
Engle, R.~F. (1984).
Wald, likelihood ratio, and Lagrange multiplier tests in econometrics.
En Z. Griliches \& M.~D. Intriligator (Eds.),
\textit{Handbook of Econometrics}, Vol.~2 (pp.\ 775--826).
North-Holland.

% ==========================================================
\end{document}
